\subsection{Novelty (Generated-Content) Probe}
\label{app:novelty_probe}

The style, implausibility, and relocation instruments all penalise things the
model does wrong. None of them catch the opposite failure: a model that
generates nothing new and simply coasts on the visible scene passes every
damage check while doing no world-modelling at all. We therefore add a novelty
probe that asks whether the model introduced genuinely new content, and we use
it as a diagnostic rather than a scored column, because generating a well-formed
new object is not itself an error.

The probe is a third binary question put to the same judge (Qwen3-VL-8B) under
the identical protocol as the style and implausibility probes
(Section~\ref{app:vlm_quality_probes}): four real REFERENCE frames, sixteen
GENERATED frames, forced answer prefix, score read as the logit probability of
\emph{Yes}. The question is:

\begin{quote}
\small
\emph{Does any NEW OBJECT appear in the generated frames that was not visible in
the reference---a vehicle, person, structure or large item that the model
introduced? Newly revealed parts of the existing scene (from driving) do NOT
count; only genuinely new objects. Answer Yes or No.}
\end{quote}

We do not report novelty as a per-model rate, because on its own it does not
separate good models from bad ones---a high rate can mean useful generation or
confabulation. Its value is that it distinguishes the two when read against the
commanded action. A model that has learnt to condition on the action should
generate new content only when the command reveals unseen scene: almost never
when driving forward along visible road, and often when turning or reversing
into geography that was never in view. A model that hallucinates instead
produces new objects regardless of direction.

Table~\ref{tab:novelty_direction} shows this split. Our model's novelty tracks
the action: near zero on forward commands, rising to $30$--$50\%$ on lateral and
backward commands where new scene must genuinely be synthesised. minWM fires at
$48$--$72\%$ in every direction, including forward, where nothing new should
appear, and it continues to introduce objects even in rollouts we classify as
static ($36\%$ novelty when frozen versus $60\%$ when moving). This is the
signature of confabulation, and it is the evidence that minWM's otherwise clean
damage scores are bought by inventing content and by not moving, rather than by
modelling the scene. At the other extreme, \texttt{noadaln} fires at only $4\%$:
it barely generates anything, consistent with its high static rate and weak
action response.

\begin{table}[t]
\centering
\small
\caption{
    Novelty-probe fire rate (\% of active rollouts, $P(\text{yes})>0.5$) by
    commanded direction, for our full-capacity model against minWM. Our model
    generates new content only where the action reveals unseen scene; minWM
    introduces objects in every direction, including forward.
}
\label{tab:novelty_direction}
\begin{tabular}{lcccccccc}
\toprule
Model & F & FL & FR & L & R & B & BL & BR \\
\midrule
Ours (pca8) & $0$  & $0$  & $13$ & $50$ & $27$ & $33$ & $47$ & $37$ \\
minWM       & $53$ & $56$ & $64$ & $48$ & $65$ & $59$ & $72$ & $61$ \\
\bottomrule
\end{tabular}
\end{table}

The probe has one measured limitation: it detects objects \emph{inserted} into a
stable scene, and is largely blind to novelty introduced by wholesale
relocation. When a model abandons the scene for a different street, the entire
frame is new content, yet the probe often reads near zero because there is no
discrete new object against a stable background. Novelty is therefore
anti-correlated with relocation and must be read jointly with the relocation
flag, never alone; this is the second reason we keep it as a diagnostic rather
than a headline column.
